\documentclass[12pt]{article}

\usepackage[a4paper, margin=0.8in]{geometry}
\usepackage{graphicx} % Required for inserting images
\usepackage{amsmath}
\usepackage{setspace}
\setlength{\parindent}{0pt}

\begin{document}

\section*{\textbf{BdOI\hspace{0.5em}Monthly\hspace{0.5em}Round\hspace{0.5em}1\hspace{0.5em}Solutions:\hspace{0.5em}Problem\hspace{0.5em}Groupchat}}

\vspace{0.3cm}

\textbf{\Large Problem Information}

\vspace{0.3cm}

{\large Problem Author: Aritro Sarkar}

\vspace{0.3cm}

{\large Problem Preparation: Aritro Sarkar, Nabeul Islam}

\vspace{0.3cm}

{\large Editorial: Aritro Sarkar, Md Sazid Hasan}

\vspace{0.3cm}



\vspace{0.5cm}

\textbf{\Large Subtask by subtask solutions}

\vspace{0.3cm}

\textbf{\large Preliminary Observation:}

\vspace{0.3cm}

\begin{spacing}{1.5}
{\large
Let $S_i$ be the last message seen by the $i$-th person. Let $m$ be the current total number of messages sent.}

{\large Let a simple simulation of events be:}


{\large
For every query of type $1$, the number of sent messages increase by $1$. Formally, $m:=m+1$.
}


{\large For every query of type $2$, friend $f$ reads till the very last message $m$. Formally, $S_f := m$.}


{\large For every query of type $3$, we are tasked with finding the number of indices $i$, such that, $S_i \ge s$.}


{\large Since $N \cdot Q \le 2.5 \cdot 10^{11}$, it is clear that a standard simulation of the events will result in a Time Limit Exceeded error. However, it will yield partial scores.}

\vspace{0.3cm}

\textbf{\large Subtask 1: $N = 1$}

\vspace{0.3cm}

{\large This subtask is quite easy. A standard simulation of events will suffice. However, there is an even easier solution. For each query of type 1, $m$, the total number of messages sent so far will increase by $1$. Formally, $m := m + 1$. For every query of type 2, we will set $S_1$ to $m$. Formally, $S_1 := m$. For every query of type 3,}

\end{spacing}

\vspace{-0.8cm}

{\large 
\[
\textbf{The answer A} =
\begin{cases}
1 & \text{if } S_1 \ge s,\\
0 & \text{otherwise} 
\end{cases}
\]}

{\large Maintaining this simple simulation with a loop will suffice.}

\vspace{0.6cm}

\newpage

\textbf{\large Subtask 2: $N,Q \le 5000$}

\vspace{0.1cm}
\begin{spacing}{1.5}
{\large For this subtask, a simulation of the events that occur will yield points. The simulation goes as follows:}
\end{spacing}
\vspace{0.1cm}

{\large Type 1: $m := m + 1$.}

\vspace{0.1cm}

{\large Type 2: $S_f := m$}

\vspace{0.1cm}

{\large Type 3:

\[
\textbf{The answer A = }
\sum_{i=1}^{n} 
\[
\begin{cases}
1 & \text{if } S_1 \ge s,\\
0 & \text{otherwise} 
\end{cases}
\]

\]
}

\vspace{0.1cm}





\textbf{\large Subtask 3: All type 3 queries appear after all other queries}

\begin{spacing}{1.5}
\vspace{0.3cm}

{\large This subtask is special. Once all the type 1 and type 2 queries have taken place, only after, we will have to answer type 3 queries and $S$ will not change again. Each type 3 queries ask us how many elements of $S$ are larger or equal to the input $s$. We can solve this sub-problem using an ordered set, but an even easier solution is using difference array and prefix sum techniques. Here's how: We will do everything as the same for types 1 and 2 queries. Let $D$ be the difference array. $D_i$ is going to help us calculate how many people have seen the $i$-th message. Initially, $D_i = 0$ for all indices. For every person from $1$ to $n$, $D_{1,S_i}$ should be increased by 1. For every person $i$, $D_i := D_i + 1$ and $D_{S_i+1} := D_{S_i+1} - 1$.}


{\large Afterwards, we will perform prefix sum on $D$.}

\vspace{0.1cm}

{\large Formally, let 
$P_i = \sum_{j=1}^{i} D_j$.
Now, $P_i$ is the number of people who have read the $i$-th message. Our answer for a type 3 query is $P_s$.}

\vspace{0.5cm}

\textbf{\Large Full solution description}

\vspace{0.3cm}

{\large Looking back at the standard simulation, the only reason it won't give full scores is that because it's time complexity is $O(N \cdot Q)$. Our simulation does type 1 and type 2 queries in $O(1)$, but the type 3 queries take $O(N)$ complexity. An ordered set or a segment tree can fully solve the problem. details are as follows:}

\newpage

{\large \textbf{Ordered Set:} For every type 1 query, we assign it a timestamp which is equal to the number of current queries and insert it to a vector which contains the assigned timestamp of every sent message, let it be $M$. For every type 2 query, if $S_f \neq 0$, then we erase $S_f$ from the ordered set. After, we update $S_f$. Now, we insert $S_f$ to the ordered set.}

\vspace{0.1cm}

{\large For every type 3 query, the answer is \textbf{os.size()} $-$ \textbf{os.order\_of\_key}($M_{s-1}$).}

\vspace{.2cm}

{\large \textbf{Segment Tree or Fenwick Tree:} Let each of the tree indices denote the timestamp. Now, similarly to the ordered set solution, when we are given a type 2 query if $S_f \neq 0$, we decrease 1 from the $S_fth$ index of the tree, and after updating $S_f$, we increase the new $S_fth$ index by 1. 
\\Now, for query of type 3, the answer is the sum of the range \textbf{[M, R]}, where R denotes the last index of the tree. }

\vspace{0.6cm}

{\large Both of these solution achieves a Time Complexity of $O(Q \cdot \log_{2}{n})$, which will suffice for this problem's constraints.  }

\end{spacing}

\end{document}
